% Автоматаар үүсгэсэн — эх файл: docs/resource-budget.md
\chapter{Нөөцийн төсөв}

Энэ бол \textbf{чиглүүлэх тооцоо}. Бодит утгыг Лаб 1-д өөрсдөө хэмжинэ.

\hypertarget{pi-raspberry-pi-3b-1-gb-ux44bux43d-ux430ux440ux438ux444ux43cux435ux442ux438ux43a}{%
\section{1. {[}Pi{]} Raspberry Pi 3B --- 1 GB-ын арифметик}\label{pi-raspberry-pi-3b-1-gb-ux44bux43d-ux430ux440ux438ux444ux43cux435ux442ux438ux43a}}

\hypertarget{ux44eux443-ux445ux430ux430ux43dux430-ux44fux432ux434ux430ux433-ux432ux44d}{%
\subsection{Юу хаана явдаг вэ}\label{ux44eux443-ux445ux430ux430ux43dux430-ux44fux432ux434ux430ux433-ux432ux44d}}

\begingroup\scriptsize\begin{longtable}[]{@{}
  >{\raggedright\arraybackslash}p{(\columnwidth - 4\tabcolsep) * \real{0.3333}}
  >{\raggedright\arraybackslash}p{(\columnwidth - 4\tabcolsep) * \real{0.3333}}
  >{\raggedright\arraybackslash}p{(\columnwidth - 4\tabcolsep) * \real{0.3333}}@{}}
\toprule
\begin{minipage}[b]{\linewidth}\raggedright
Хэсэг
\end{minipage} & \begin{minipage}[b]{\linewidth}\raggedright
MiB
\end{minipage} & \begin{minipage}[b]{\linewidth}\raggedright
Тэмдэглэл
\end{minipage} \\
\midrule
\endhead
\bottomrule
\endlastfoot
Нийт физик санах ой & 1024 & \\
GPU-д (\texttt{gpu\_mem=16}) & −16 & анхдагч 64 → 48 MiB хэмнэнэ \\
Цөм + firmware & \textasciitilde85 & \\
\textbf{Системд харагдах (\texttt{MemTotal})} & \textbf{\textasciitilde925} & \texttt{free\ -m} энийг харуулна \\
Pi OS Lite, headless, амарч байгаа & −90\ldots−130 & systemd, sshd, journald \\
Docker демон & −60\ldots−80 & \texttt{dockerd} + \texttt{containerd} \\
\textbf{Ажилд боломжтой} & \textbf{\textasciitilde720\ldots780} & \\
\end{longtable}\endgroup

\hypertarget{ux438ux440ux43cux44dux433ux438ux439ux43d-ux4afux439ux43bux447ux438ux43bux433ux44dux44dux43dux438ux439-ux442ux4e9ux441ux4e9ux432}{%
\subsection{Ирмэгийн үйлчилгээний төсөв}\label{ux438ux440ux43cux44dux433ux438ux439ux43d-ux4afux439ux43bux447ux438ux43bux433ux44dux44dux43dux438ux439-ux442ux4e9ux441ux4e9ux432}}

\begingroup\scriptsize\begin{longtable}[]{@{}
  >{\raggedright\arraybackslash}p{(\columnwidth - 6\tabcolsep) * \real{0.2500}}
  >{\raggedright\arraybackslash}p{(\columnwidth - 6\tabcolsep) * \real{0.2500}}
  >{\raggedright\arraybackslash}p{(\columnwidth - 6\tabcolsep) * \real{0.2500}}
  >{\raggedright\arraybackslash}p{(\columnwidth - 6\tabcolsep) * \real{0.2500}}@{}}
\toprule
\begin{minipage}[b]{\linewidth}\raggedright
Үйлчилгээ
\end{minipage} & \begin{minipage}[b]{\linewidth}\raggedright
Compose хязгаар
\end{minipage} & \begin{minipage}[b]{\linewidth}\raggedright
Бодит амралтын хэрэглээ
\end{minipage} & \begin{minipage}[b]{\linewidth}\raggedright
Тэмдэглэл
\end{minipage} \\
\midrule
\endhead
\bottomrule
\endlastfoot
\texttt{mosquitto} & 128 MiB & 12--25 MiB & Дараалал дүүрэхэд өснө: 100k мессеж × 250 B ≈ 25 MiB \\
ирмэгийн агент (venv) & 200 MiB (systemd) & 45--70 MiB & numpy \textasciitilde25 MiB эзэлнэ \\
+ TFLite дүгнэлт (Лаб 6) & ↑ & +40\ldots+90 MiB & загварын хэмжээ ба цонхноос хамаарна \\
VS Code Remote сервер & --- & \textbf{150--250 MiB} & ⚠ хэмжилтийн өмнө таслах! \\
\end{longtable}\endgroup

\textbf{Нийлбэр (Лаб 1--5):} \textasciitilde120 MiB → \textbf{600 MiB илүүдэлтэй.} Тав тухтай.
\textbf{Нийлбэр (Лаб 6, дүгнэлттэй):} \textasciitilde200 MiB → \textbf{520 MiB илүүдэлтэй.} Хангалттай.
\textbf{Лаб 8 (K3s):} доорх §3-ыг үзнэ үү --- энэ нь хамгийн чанга.

\hypertarget{ux445ux44dux437ux44dux44d-ux430ux43dux445ux430ux430ux440ux430ux445-ux432ux44d}{%
\subsection{Хэзээ анхаарах вэ}\label{ux445ux44dux437ux44dux44d-ux430ux43dux445ux430ux430ux440ux430ux445-ux432ux44d}}

\begin{Shaded}
\begin{Highlighting}[]
\FunctionTok{free} \AttributeTok{{-}m} \KeywordTok{|} \FunctionTok{awk} \StringTok{\textquotesingle{}/Mem:/\{print $7 " MiB available"\}\textquotesingle{}}
\end{Highlighting}
\end{Shaded}

\begingroup\scriptsize\begin{longtable}[]{@{}
  >{\raggedright\arraybackslash}p{(\columnwidth - 2\tabcolsep) * \real{0.5000}}
  >{\raggedright\arraybackslash}p{(\columnwidth - 2\tabcolsep) * \real{0.5000}}@{}}
\toprule
\begin{minipage}[b]{\linewidth}\raggedright
Үлдэгдэл
\end{minipage} & \begin{minipage}[b]{\linewidth}\raggedright
Утга
\end{minipage} \\
\midrule
\endhead
\bottomrule
\endlastfoot
\textgreater{} 400 MiB & хэвийн \\
150--400 MiB & болгоомжтой --- шинэ үйлчилгээ нэмэхээс өмнө хэмж \\
\textless{} 150 MiB & \textbf{аюултай} --- swap идэвхжиж, бүх саатал 10--100 дахин өснө \\
swap \texttt{si/so} \textgreater{} 0 & аль хэдийн хязгаарт хүрсэн. Энэ нь \textbf{хэмжилт}, алдаа биш. \\
\end{longtable}\endgroup

\texttt{tools/\allowbreak{}measure\_stack.\allowbreak{}sh\ -\/\allowbreak{}-role\ edge} эдгээрийг автоматаар шалгана.

\hypertarget{ux437ux43a-ux437ux4e9ux4e9ux432ux440ux438ux439ux43d-ux43aux43eux43cux43fux44cux44eux442ux435ux440-ux4afux4afux43bux43dux438ux439-ux434ux430ux432ux445ux430ux440ux433ux430}{%
\section{2. {[}ЗК{]} Зөөврийн компьютер --- үүлний давхарга}\label{ux437ux43a-ux437ux4e9ux4e9ux432ux440ux438ux439ux43d-ux43aux43eux43cux43fux44cux44eux442ux435ux440-ux4afux4afux43bux43dux438ux439-ux434ux430ux432ux445ux430ux440ux433ux430}}

\begingroup\scriptsize\begin{longtable}[]{@{}
  >{\raggedright\arraybackslash}p{(\columnwidth - 8\tabcolsep) * \real{0.2000}}
  >{\raggedright\arraybackslash}p{(\columnwidth - 8\tabcolsep) * \real{0.2000}}
  >{\raggedright\arraybackslash}p{(\columnwidth - 8\tabcolsep) * \real{0.2000}}
  >{\raggedright\arraybackslash}p{(\columnwidth - 8\tabcolsep) * \real{0.2000}}
  >{\raggedright\arraybackslash}p{(\columnwidth - 8\tabcolsep) * \real{0.2000}}@{}}
\toprule
\begin{minipage}[b]{\linewidth}\raggedright
Үйлчилгээ
\end{minipage} & \begin{minipage}[b]{\linewidth}\raggedright
Профайл
\end{minipage} & \begin{minipage}[b]{\linewidth}\raggedright
Compose хязгаар
\end{minipage} & \begin{minipage}[b]{\linewidth}\raggedright
Хүлээгдэх хэрэглээ
\end{minipage} & \begin{minipage}[b]{\linewidth}\raggedright
Тэмдэглэл
\end{minipage} \\
\midrule
\endhead
\bottomrule
\endlastfoot
\texttt{emqx} & core & 1024 MiB & 150--250 MiB & Холболт нэмэгдэхэд \textasciitilde4 KiB/холболт \\
\texttt{influxdb} & core & 1024 MiB & 250--500 MiB & Бичилтийн ачаалалтай өснө \\
\texttt{grafana} & core & 512 MiB & 120--200 MiB & Самбарын тоогоор өснө \\
\texttt{registry} & core & 384 MiB & 60--110 MiB & FastAPI + SQLite \\
\texttt{nodered} & pipeline & 512 MiB & 120--200 MiB & Урсгалын нарийвчлалаас \\
\texttt{dex} & app & 256 MiB & 30--60 MiB & \\
\texttt{graphql-api} & app & 512 MiB & 80--150 MiB & Python/FastAPI \\
\texttt{ollama} & ai & 4096 MiB & загвараас & qwen2.5:1.5b q4 ≈ 1.2--1.6 GiB + KV \\
\texttt{thingsboard} & tb (сонголт) & 2500 MiB & 1200--1800 MiB & зөвхөн Лаб 2-ын харьцуулалт \\
\end{longtable}\endgroup

\hypertarget{ux43fux440ux43eux444ux430ux439ux43bux44bux43d-ux43dux438ux439ux43bux431ux44dux440}{%
\subsection{Профайлын нийлбэр}\label{ux43fux440ux43eux444ux430ux439ux43bux44bux43d-ux43dux438ux439ux43bux431ux44dux440}}

\begingroup\scriptsize\begin{longtable}[]{@{}llll@{}}
\toprule
Хослол & Хязгаарын нийлбэр & Практикт & Docker Desktop-д хэрэгтэй \\
\midrule
\endhead
\bottomrule
\endlastfoot
core & 2.9 GiB & 0.6--1.1 GiB & 4 GB \\
core + pipeline & 3.4 GiB & 0.7--1.3 GiB & 4 GB \\
core + pipeline + app & 4.2 GiB & 0.8--1.5 GiB & 4 GB \\
+ ai (Лаб 8) & 8.2 GiB & 2.0--3.1 GiB & \textbf{6 GB} \\
core + tb (Лаб 2 сонголт) & 5.4 GiB & 1.8--2.9 GiB & \textbf{6 GB} \\
\end{longtable}\endgroup

\begin{quote}
\textbf{Харьцуулалт:} ThingsBoard ГАНЦААРАА 1.2--1.8 GiB эзэлнэ --- энэ нь Pi 3B-ийн НИЙТ санах ойноос \textbf{хоёр дахин их}. Лаб 2-ын хяналтын асуултад: "яагаад ирмэг дээр бүрэн платформ ажиллуулж болохгүй вэ?" гэдгийн тоон хариулт нь энэ.
\end{quote}

\hypertarget{pi-ux43bux430ux431-8-k3s-ux438ux439ux43d-ux441ux430ux43dux430ux445-ux43eux439ux43d-ux442ux4e9ux441ux4e9ux432}{%
\section{3. {[}Pi{]} Лаб 8 --- K3s-ийн санах ойн төсөв}\label{pi-ux43bux430ux431-8-k3s-ux438ux439ux43d-ux441ux430ux43dux430ux445-ux43eux439ux43d-ux442ux4e9ux441ux4e9ux432}}

Лаб 8-д Pi дээр \textbf{Docker-ийг зогсоож}, K3s-ийг ажиллуулна. Хоёуланг зэрэг ажиллуулах орон зай байхгүй.

\begingroup\scriptsize\begin{longtable}[]{@{}lll@{}}
\toprule
Хэсэг & Requests & Бодит \\
\midrule
\endhead
\bottomrule
\endlastfoot
k3s server (containerd + kubelet + API + sqlite) & --- & 250--350 MiB \\
coredns & 70 MiB & \textasciitilde25 MiB \\
metrics-server (HPA-д ЗААВАЛ) & 70 MiB & \textasciitilde30 MiB \\
local-path-provisioner & --- & \textasciitilde10 MiB \\
\texttt{mosquitto} pod & 64 MiB & 12--25 MiB \\
\texttt{edge-agent} pod & 64 MiB & 45--70 MiB \\
\textbf{Нийт requests} & \textbf{268 MiB} & \\
\textbf{Бодит нийлбэр} & & \textbf{\textasciitilde400--510 MiB} \\
Allocatable (925 − цөм − k3s үндэс) & \textasciitilde825 MiB & \\
\end{longtable}\endgroup

\textbf{Дүгнэлт:} ажиллана, гэхдээ \texttt{traefik} ба \texttt{servicelb}-ийг заавал унтраана (тус бүр 60--90 MiB). HPA 2 хувь болтол өргөжиж чадна; 3 дахь хувь OOMKilled болно --- \textbf{энэ нь Лаб 8-ын гол ажиглалт}.

\begin{Shaded}
\begin{Highlighting}[]
\ExtensionTok{curl} \AttributeTok{{-}sfL}\NormalTok{ https://get.k3s.io }\KeywordTok{|} \FunctionTok{sh} \AttributeTok{{-}s} \AttributeTok{{-}} \DataTypeTok{\textbackslash{}}
  \AttributeTok{{-}{-}disable}\NormalTok{ traefik }\AttributeTok{{-}{-}disable}\NormalTok{ servicelb }\AttributeTok{{-}{-}write{-}kubeconfig{-}mode}\NormalTok{ 644}
\end{Highlighting}
\end{Shaded}

\hypertarget{ux441ux430ux43dux430ux43cux436}{%
\section{4. Санамж}\label{ux441ux430ux43dux430ux43cux436}}

\begin{itemize}
\tightlist
\item
  \texttt{deploy.resources.limits} нь \textbf{хязгаар} л тавьдаг, нөөц \textbf{баталгаажуулдаггүй}. Pi дээр хязгаарыг давсан контейнерийг цөм OOM-оор устгана; лог нь \texttt{docker\ inspect\ \textless{}c\textgreater{}\ \textbar{}\ grep\ OOMKilled}.
\item
  \texttt{docker\ stats}-ын санах ойн тоо зөв гарахад \texttt{cgroup\_enable=memory} хэрэгтэй (SETUP.md Б.5).
\item
  Идэвхтэй хөргөлт байхгүй бол Pi 3B ачаалал дор 80 °C давж, давтамжаа бууруулна --- бүх хэмжилт гажина. \texttt{vcgencmd\ get\_throttled} нь \textbf{\texttt{0x0}} байх ёстой.
\item
  microSD-гийн бичих хурд (10--25 MB/s) нь InfluxDB, mosquitto persistence, swap гурвыг зэрэг хязгаарлана. Лаб 5-д үүнийг бодитоор хэмжинэ.
\item
  Хэмжилтийн өмнө \textbf{VS Code Remote серверийг таслах}: \texttt{pkill\ -f\ vscode-server}. 150--250 MiB ялгаа гарна.
\end{itemize}
